\chapter{代数拓扑}\label{appx:Alge-Topo}

- \(H^l_{\mathrm{dR}}(U)\)：区域 \(U\) 的 **\(l\) 阶 de Rham 上同调**，即“闭 \(l\)-形式模去恰当 \(l\)-形式”。它为零，意味着每个闭 \(l\)-形式都有全局势。
- \(H^2(X;\mathbb R)\)：空间 \(X\) 的**实系数二阶上同调**。若 \(X\) 是光滑流形，de Rham 定理给出
  \[
  H^2(X;\mathbb R)\cong H^2_{\mathrm{dR}}(X).
  \]

上标 \(2\) 表示**二阶**，不是“第二种”。讨论一般 \(l\)-形式时应写 \(l\)；只有讨论闭 \(2\)-形式时才对应 \(H^2\)。